% =============================================================================
% Section 2: Core Finding — The Defence–Logistics Duplex
% =============================================================================
\section{Core Finding: The Defence--Logistics Duplex}%
\label{sec:core-finding}

This section presents the central empirical result of the paper in
three parts: the variance decomposition, the axis dominance pattern,
and the rank-disruption asymmetry.

\subsection{Variance share: 69.4\% from two axes}

The {\ISI} composite is an unweighted arithmetic mean of six axes
(\cref{eq:composite}).  Under this aggregation rule, the population
variance of the composite decomposes exactly into contributions from
each axis's own variance and from all pairwise covariances
(\cref{sec:proof}).

\Cref{tab:variance-decomp} reports the result.  Defence contributes
35.10\% of total composite variance; logistics contributes 34.26\%.
The combined share is 69.37\%.  The remaining four axes
contribute:

\begin{itemize}[leftmargin=2em]
  \item Critical inputs: 20.92\%
  \item Technology: 5.23\%
  \item Energy: 3.46\%
  \item Financial: 1.03\%
\end{itemize}

\noindent
The own-variance component (the contribution of each axis's marginal
variance, absent covariance) accounts for 64.28\% of total composite
variance.  Cross-covariances account for the remaining 35.72\%.

Two features of this decomposition warrant emphasis:

\begin{enumerate}[leftmargin=2em]
  \item \textbf{The duplex is nearly symmetric.}  Defence and
    logistics contribute almost equal shares (35.1\% vs.\ 34.3\%).
    Neither axis alone dominates; both are necessary to explain the
    composite's cross-country differentiation.

  \item \textbf{Financial is negligible.}  The financial axis
    contributes 1.03\% of composite variance.  Cross-border banking
    and portfolio dependencies, while economically significant, do not
    differentiate EU-27 countries from one another under the {\HHI}
    concentration metric.
\end{enumerate}

\subsection{Axis dominance: 16/27 and 11/27}

The variance decomposition describes the distribution of
\emph{dispersion}.  A complementary question is: which axis scores
highest for each country?

\Cref{tab:dominance-counts} reports the result.  In 16 of
27~countries (59.3\%), the defence axis is the highest-scoring axis.
In the remaining 11~countries (40.7\%), the logistics axis is
highest.  No country is dominated by financial, energy, technology,
or critical inputs.

This is a strong structural pattern.  It means that every EU-27
country's sovereignty-risk profile is, at the country level,
\emph{either} a defence story \emph{or} a logistics story.  The
other four axes contribute to the composite score but never determine
a country's peak vulnerability.

Among the 16~defence-dominant countries, the mean gap between the
defence score and the second-highest axis is 0.184 on the unit
interval.  Among the 11~logistics-dominant countries, the
corresponding gap is 0.088.  Defence dominance is not only more
prevalent but more pronounced.

\subsection{Rank disruption: defence removal is most destructive}

Leave-one-axis-out (LOO) analysis removes each axis in turn and
recomputes the composite as the mean of the remaining five.  The
resulting ranking is compared to the baseline using Spearman~$\rho$
and Kendall~$\tau$.

\Cref{tab:loo-disruption} reports the result.  Removing the defence
axis produces the lowest rank-order correlation with the baseline:
$\rho = 0.833$, $\tau = 0.675$.  Removing logistics produces
$\rho = 0.926$, $\tau = 0.789$.  By contrast, removing financial or
energy barely perturbs the ranking ($\rho > 0.99$).

The asymmetry between defence and logistics in the LOO test
(despite near-equal variance shares) reflects a structural
difference: defence axis scores are more dispersed (population
standard deviation 0.244 vs.\ 0.205 for logistics) and less
correlated with other axes, so their removal produces larger rank
displacements.

The maximum absolute rank displacement when defence is removed is
10~positions.  The mean absolute displacement is 3.33~positions.
For financial removal, the mean displacement is 0.67~positions.

\subsection{Summary: the duplex}

The three results---variance share, axis dominance, and
rank disruption---converge on a single structural fact:

\begin{tcolorbox}[
  colback=ISILight,
  colframe=ISIBlue,
  fonttitle=\bfseries\sffamily,
  title={The defence--logistics duplex}
]
The cross-country dispersion of the {\ISI} composite in the EU-27 is
structurally determined by two axes.  Defence and logistics together
account for 69.4\% of composite variance, dominate the axis profile of
every country, and produce the largest rank disruptions when omitted.
The remaining four axes are second-order for cross-country
differentiation.
\end{tcolorbox}
